% Preamble
\documentclass[11pt, aspectratio=169]{beamer}
	% Packages
	\usepackage[english]{babel}
	\usepackage{lipsum}
	\usepackage{mathptmx}
	\usepackage{xcolor}
	% Themes
	\usefonttheme[onlymath]{serif}
	\usetheme[nofirafonts]{focus}
	\renewcommand{\thempfootnote}{\arabic{mpfootnote}}
	% Cover
	\title{Microeconometrics Using Stata}
	\subtitle{Linear regression basics: Exercises}
	\author{
		Jhon R. Ordoñez
		\footnote{\label{1} National University of San Cristóbal de Huamanga}
		\footnote{\label{2} Faculty of Economic, Administrative and Accounting Sciences}
		\footnote{\label{3} Professional School of Economics}
	}
	\date{\today}
% Body
\begin{document}
	% Cover
	\begin{frame}
		\maketitle
	\end{frame}
	% Outline
	\begin{frame}[t]
		\frametitle{Outline}
		\tableofcontents
	\end{frame}
	% Section 1
	\section{Exercise 1}
		% Slide 1.1
		\begin{frame}
			\frametitle{Exercise 1}
				\begin{block}{Exercise 1}
					Fit the model in \textcolor{red!50}{\textbf{section 3.5}} using only the first $100$ observations.
					Compute standard errors in three ways: default, heteroskedastic, and
					cluster-robust where clustering is on the number of chronic problems.
					Use \textcolor{blue}{\textbf{estimates}} to produce a table with three sets of coefficients and
					standard errors, and comment on any appreciable differences in the
					standard errors. Construct a similar table for three alternative sets of
					heteroskedasticity-robust standard errors, obtained by using the
					\textcolor{blue}{\textbf{vce(robust)}}, \textcolor{blue}{\textbf{vce(hc2)}}, and \textcolor{blue}{\textbf{vce(hc3)}} options, and comment on any
					differences between the different estimates of the standard errors.
				\end{block}
		\end{frame}
	% Section 2
	\section{Exercise 2}
		% Slide 2.1
		\begin{frame}
			\frametitle{Exercise 2}
			\begin{block}{Exercise 2}
				Fit the model in \textcolor{red!50}{\textbf{section 3.5}} with robust standard errors reported. Test
				at $5\%$ the joint significance of the demographic variables \textcolor{blue}{\textbf{age}}, \textcolor{blue}{\textbf{female}},
				and \textcolor{blue}{\textbf{income}}. Test the hypothesis that being male (rather than female)
				has the same impact on medical expenditures as aging 10 years. Fit the
				model under the constraint that $\beta_{phylim} = \beta_{actlim}$ by first typing
				\textcolor{blue}{\textbf{constraint 1 phylim = actlim}} and then by using \textcolor{blue}{\textbf{cnsreg}} with the
				\textcolor{blue}{\textbf{constraints(1)}} option.
			\end{block}
		\end{frame}
	% Section 3
	\section{Exercise 3}
		% Slide 3.1
		\begin{frame}
			\frametitle{Exercise 3}
			\begin{block}{Exercise 3}
				Fit the model in \textcolor{red!50}{\textbf{section 3.6}}, and implement the RESET test manually by
				regressing $y$ on $x$ and $\hat{y}^{2}$, $\hat{y}^{3}$, and $\hat{y}^{4}$ and jointly testing that the
				coefficients of $\hat{y}^{2}$, $\hat{y}^{3}$, and $\hat{y}^{4}$ are $0$. 
				To get the same results as \textcolor{blue}{\textbf{estat ovtest}}, 
				do you need to use default or robust estimates of the \textbf{VCE} in
				this regression? Comment. Similarly, implement \textcolor{blue}{\textbf{linktest}} by
				regressing $y$ on $\hat{y}$ and $\hat{y}^{2}$ and testing that the coefficient of $\hat{y}^{2}$ is $0$. To
				get the same results as \textcolor{blue}{\textbf{linktest}}, do you need to use default or robust
				estimates of the \textbf{VCE} in this regression? Comment.
			\end{block}
		\end{frame}
	% Section 4
	\section{Exercise 4}
		% Slide 4.1
		\begin{frame}
			\frametitle{Exercise 4}
			\begin{block}{Exercise 4}
				a
			\end{block}
		\end{frame}
	% References
	\begin{frame}[t,allowframebreaks]
		\frametitle{References}
		\bibliography{biblio.bib}
		% Style
		\bibliographystyle{apalike}
		% Authors
		\nocite{cameron2022-I}
		\nocite{cameron2022-II}
	\end{frame}
	% Cover
	\begin{frame}
		\maketitle
	\end{frame}
\end{document}